\section{방법론}
\label{sec:ch2}

이 장은 최종 모델 C-BAT의 구조, 추정, 예측 산출 방식을 정의한다. 자료 전처리는 3장, 검증 결과와 비교 모델은 4장에서 다룬다. 이 장에서 정한 기호는 보고서 전체에서 같은 뜻으로 쓴다.

\subsection{문제 정의와 표기}
\label{sec:ch2-problem}

예측 과제는 예측일 $d$의 00:00에 그날 15분 슬롯 전력 $y_{d,t}$, $t=1,\dots,96$을 한꺼번에 내는 것이다. 00:00에 쓸 수 있는 정보는 그날의 시간별 생산계획 $q_{d,h}$($h=0,\dots,23$), 달력 정보(요일, 날짜유형), 그리고 $d-1$일까지 관측된 전력이다. 예측일 당일의 전력은 한 슬롯도 쓰지 않는다.

모델은 점예측 하나가 아니라 하루 전체 전력 경로의 예측분포를 낸다. 이 분포에서 세 가지 산출물을 얻는다. 첫째는 슬롯별 분위수와 점예측이다. 둘째는 일 피크 $M_d=\max_t y_{d,t}$의 분포다. 셋째는 피크 초과확률 $P(M_d > C)$이며, $C$는 학습 구간에서 정한 피크 기준이다. 슬롯마다 따로 만든 분위수로는 일 피크 분포를 얻을 수 없으므로, 하루 곡선 전체를 함께 표집하는 구조가 필요하다. 표~\ref{tab:ch2-notation}은 이 장에서 쓰는 기호다.

\begin{table}[htbp]
\centering
\caption{주요 기호.}
\label{tab:ch2-notation}
\small
\begin{tabular}{ll}
\toprule
기호 & 뜻 \\
\midrule
$d$ & 날짜(예측일) \\
$t=1,\dots,96$ & 15분 슬롯 \\
$h=0,\dots,23$ & 생산계획의 시각 \\
$y_{d,t}$ & 슬롯 전력(kW) \\
$s_y$, $z_{d,t}=y_{d,t}/s_y$ & 학습 전력의 표준편차와 표준화 전력(평균은 빼지 않음) \\
$q_{d,h}$ & 시간별 생산계획량, $q_{\max}$는 학습 구간 최댓값 \\
$k\in\{0,1,2,3\}$ & 가동유형(계획 생산시간 0 / 1--12 / 13--19 / 20시간 이상) \\
$g\in\{\text{비가동},\text{가동}\}$ & 잡음 상태($k=0$이면 비가동, $k\ge1$이면 가동) \\
$\mu_{k,t}$ & 가동유형별 기본곡선(표준화 단위) \\
$r_{d,t}$, $\gamma_k$ & 표준화한 기준일 곡선과 그 계수 \\
$x_c(q)$, $c\in\{\mathrm{on},\log\}$ & 생산 채널(생산 여부, 정규화 로그 생산량) \\
$w_{c,k}=e^{\omega_{c,k}}$ & 채널별 생산 이득(양수) \\
$\alpha_{t,h}$ & 시간 계획을 슬롯으로 옮기는 커널 가중치 \\
$\ell$ & 커널 시차(시간), 최종 모델에서는 0 \\
$u_d$, $\sigma_{u,g}$ & 날짜 효과와 그 표준편차 \\
$e_{d,t}$, $\rho_g$, $\sigma_{\eta,g}$ & 슬롯 AR(1) 잡음, 자기상관, 혁신 표준편차 \\
$\omega_d$ & 반감기 시간 가중치 \\
$M_d$ & 일 피크 $\max_t y_{d,t}$ \\
$C_{50}, C_{75}, C_{90}$ & 학습 구간 가동일 피크의 50·75·90\% 분위수 \\
$B_d$ & 월 최대 경보의 기준선(00:00에 알려진 요금 적용 최대) \\
\bottomrule
\end{tabular}
\end{table}

\subsection{모델 구조}
\label{sec:ch2-model}

C-BAT는 먼저 학습 구간에서 관측된 전력의 표준편차 $s_y$를 계산하고 $z_{d,t}=y_{d,t}/s_y$로 변환한다. 평균을 빼는 중심화는 하지 않는다. 이후 기본곡선, 기준일 곡선, 날짜 효과와 슬롯 잡음은 모두 이 표준화 단위로 정의한다. $q_{\max}$도 해당 학습 구간에서만 계산하며, 검증일과 9월의 전력이나 생산량으로 척도를 갱신하지 않는다.

표준화 슬롯 전력은 평균 곡선 세 항과 변동 두 항의 합이다.
\begin{equation}
z_{d,t} \;=\; \underbrace{\mu_{k,t}}_{\text{기본곡선}}
\;+\; \underbrace{\gamma_k\, r_{d,t}}_{\text{기준일}}
\;+\; \underbrace{\sum_{c\in\{\mathrm{on},\,\log\}} w_{c,k} \sum_{h=0}^{23} \alpha_{t,h}\, x_c(q_{d,h})}_{\text{생산 전달항}}
\;+\; \underbrace{u_d}_{\text{날짜 효과}}
\;+\; \underbrace{e_{d,t}}_{\text{슬롯 잡음}}
\label{eq:ch2-cbat}
\end{equation}
여기서 $k$는 예측일 $d$의 가동유형이다. 앞의 세 항은 생산계획과 최근 운전 상태가 정하는 평균 곡선이고 뒤의 두 항은 그 주변의 변동이다. 이름의 C(conditional)는 변동의 크기와 지속성을 가동 상태별로 따로 두는 데서 붙었다. 예측할 때 식의 오른쪽 전체에 $s_y$를 곱해 kW로 되돌린다.

가동유형 $k$는 그날 생산계획만으로 정해지므로 00:00에 알 수 있다. 계획량이 양수인 시간 수가 0이면 $k=0$(비가동), 1--12시간이면 $k=1$, 13--19시간이면 $k=2$, 20시간 이상이면 $k=3$이다. 기본곡선, 기준일 계수, 생산 이득은 모두 가동유형마다 따로 둔다. 운전 시간이 다른 날은 하루 곡선의 모양 자체가 다르기 때문이다.

기본곡선 $\mu_{k,t}$는 가동유형별 하루 평균 모양을 맡는다. 96개 값에 순환 2차 확률보행(RW2) 사전분포 \citep{rue2005gmrf}를 둔다. 인접 슬롯의 이차 차분 $\mu_{k,t+1}-2\mu_{k,t}+\mu_{k,t-1}$에 벌점을 주고 슬롯 96과 슬롯 1을 이어 자정 전후도 매끄럽게 잇는다. 평활 강도 $\tau_k$도 가동유형마다 추정하므로, 자료가 적은 유형에서는 곡선이 더 매끄럽게 수축된다.

기준일 곡선 $r_{d,t}$는 최근 공장 상태를 평균에 들인다. 기준일은 예측일 이전 28일 안에서 가동 여부와 날짜유형(평일·토요일·일요일)이 모두 같은 완전 관측일(96슬롯 모두 관측) 가운데 가장 최근 날이다. 이 규칙은 기준선 B0$'$와 같다. 28일 안에 그런 날이 없으면 전체 과거에서 가동 여부만 같은 가장 최근 완전 관측일을 쓰고, 그것도 없으면 학습 구간의 해당 가동유형 평균 곡선을 쓴다. 실제 기준일 전력도 $s_y$로 나누어 $r_{d,t}$를 만든다. 계수 $\gamma_k$는 기준일 곡선을 얼마나 따를지 정한다. 기본곡선이 긴 기간의 평균 모양을 맡는다면, 기준일 항은 계절에 따라 움직이는 최근 수준을 따라간다.

생산 전달항은 시간별 생산계획을 15분 슬롯의 전력 증가분으로 옮긴다. 계획량은 두 채널로 바꾼다.
\begin{equation}
x_{\mathrm{on}}(q)=\mathbf 1[q>0],\qquad
x_{\log}(q)=\frac{\log(1+q)}{\log(1+q_{\max})}.
\label{eq:ch2-channels}
\end{equation}
첫 채널은 라인이 도는지를, 둘째 채널은 생산량의 크기를 나타낸다. 이득은 $w_{c,k}=\exp(\omega_{c,k})$로 두어 양수로 제약한다. 같은 가동유형과 기준일을 유지한 채 생산량만 늘리면 두 채널과 생산 전달항은 줄지 않는다. 생산시간 수가 달라져 가동유형이 바뀌면 기본곡선과 계수도 바뀌므로 전체 예측의 단조성까지 보장하는 것은 아니다. 비가동일($k=0$)에는 생산이 없으므로 전달항을 두지 않는다($w_{c,0}=0$).

시간 계획을 슬롯으로 옮기는 가중치 $\alpha_{t,h}$는 고정한 sparsemax 커널이다.
\begin{equation}
\alpha_{t,\cdot} \;=\; \operatorname{sparsemax}\!\left(-\frac{\lvert \Delta_{t,h} - \ell\rvert^{\beta}}{\sigma}\right)_{h=0,\dots,23},
\qquad \Delta_{t,h}=\frac{t-4h-2.5}{4},
\qquad (\sigma,\beta,\ell)=(1.5,\ 2,\ 0).
\label{eq:ch2-kernel}
\end{equation}
$\Delta_{t,h}$는 슬롯 $t$의 중앙 시각과 시각 $h$의 중앙 시각 사이의 차이(시간 단위)다. sparsemax \citep{martins2016sparsemax}는 점수 벡터 $v$를 확률 단체에 유클리드 사영한 $\operatorname{sparsemax}(v)_h=\max(v_h-\tau(v),0)$이다. 가중치의 합은 1이고 먼 시각의 가중치는 정확히 0이 된다. 척도 $\sigma=1.5$, 형상 $\beta=2$, 시차 $\ell=0$은 사전분포의 중심값이며 검증 점수를 보고 고르지 않았다. $\sigma$는 이 점수 함수의 척도이며 가우시안 잡음의 표준편차가 아니다. 이 값에서 각 슬롯은 인접한 한두 개 시각에만 가중치를 둔다. 첫 슬롯과 마지막 슬롯은 한 시각, 나머지 94개 슬롯은 두 시각을 쓴다. 예컨대 00:30--00:45 슬롯은 0시 계획에 0.75, 1시 계획에 0.25를 둔다. 결과적으로 이 커널은 인접한 두 시각 계획을 시차 없이 선형 보간하는 국소 전달함수이며 분포시차 전달함수 \citep{welty2009bdlm}의 가장 단순한 형태다.

날짜 효과와 슬롯 잡음은 잡음 상태 $g$마다 따로 둔다. $g$는 비가동($k=0$)과 가동($k\ge1$)의 두 상태다.
\begin{equation}
u_d \sim \mathcal N\!\big(0,\ \sigma^2_{u,g}\big),
\qquad
e_{d,t} = \rho_g\, e_{d,t-1} + \eta_{d,t},\quad \eta_{d,t}\sim \mathcal N\!\big(0,\ \sigma^2_{\eta,g}\big).
\label{eq:ch2-noise}
\end{equation}
날짜 효과 $u_d$는 하루 전체를 함께 올리거나 내리는 변동이고, $e_{d,t}$는 슬롯 사이에 이어지는 1차 자기회귀 잡음이다. 하루의 첫 슬롯은 정상분포 $e_{d,1}\sim\mathcal N\big(0,\sigma^2_{\eta,g}/(1-\rho_g^2)\big)$에서 시작한다. 조용한 비가동일과 생산 부하가 오르내리는 가동일은 변동의 크기와 지속 시간이 다르다. 두 상태가 같은 잡음을 공유하면 비가동일의 경로가 가동일 수준으로 흔들려 일 피크가 과대평가된다. 상태별 잡음은 이 문제를 겨냥한다.

우도는 관측된 슬롯만으로 계산하고 결측 슬롯은 대체하지 않는다. 계측 결손으로 같은 날 관측 슬롯 $t$와 바로 앞 관측 슬롯 $t-s$ 사이에 $s-1$개 슬롯이 비어 있으면, AR(1)의 $s$단계 전이분포를 쓴다.
\begin{equation}
e_{d,t} \mid e_{d,t-s} \;\sim\; \mathcal N\!\left(\rho_g^{\,s}\, e_{d,t-s},\ \ \sigma^2_{\eta,g}\,\frac{1-\rho_g^{\,2s}}{1-\rho_g^{\,2}}\right).
\label{eq:ch2-gap}
\end{equation}
그날의 첫 관측 슬롯은 정상분포에서 나온다고 둔다. 이렇게 하면 결측 구간이 잡음 추정을 왜곡하지 않고 결측이 없는 날에는 식~\eqref{eq:ch2-noise}의 우도와 같아진다. 관측오차는 가우시안이다.

\subsection{베이지안 추정}
\label{sec:ch2-estimation}

모수는 시간 가중 우도를 쓰는 Gibbs--Metropolis 표집으로 추정한다. 날짜 $d$의 관측은 우도에서 가중치
\begin{equation}
\omega_d = 2^{-(d_{\max}-d)/30}
\label{eq:ch2-weight}
\end{equation}
를 받는다. $d_{\max}$는 전력이 하나 이상 관측된 학습일의 마지막 날이고 30일 전 자료의 가중치는 절반이다. 공장 전력 수준의 최근 변화를 반영하기 위해 다음 시간 가중 사후분포를 사용한다.
\begin{equation}
\pi(\Theta,\boldsymbol u\mid\boldsymbol z)
\;\propto\;\pi(\Theta)\prod_{d\in\mathcal T}
p(u_d\mid\Theta)\,
p(\boldsymbol z_{d,\mathcal O_d}\mid\Theta,u_d)^{\omega_d}.
\label{eq:ch2-weighted-posterior}
\end{equation}
$\mathcal T$는 유효 학습일, $\mathcal O_d$는 그날 관측 슬롯, $\Theta$는 날짜 효과를 제외한 모수다. 가중치는 AR(1) 우도 전체에 지수로 적용한다. 따라서 가우시안 이차항의 정밀도뿐 아니라 우도의 정규화항도 가중하며 날짜 효과의 사전분포에는 곱하지 않는다. 반감기 30일은 최종 설정이며, 이 분포는 모든 날짜를 같은 무게로 둔 사후분포와 구별한다.

평균 모수 $(\mu_{k,\cdot},\gamma_k)$는 가동유형마다 결합 정규 조건부분포에서 한꺼번에 뽑는다. 현재 $\rho_g$로 AR(1) 잔차를 백색화하면 우도는 $(\mu_{k,\cdot},\gamma_k)$에 대한 가중 최소제곱 꼴이 된다. 그 정규방정식에 RW2 벌점 $\mathbf R/\tau_k$와 $\gamma_k$의 사전 정밀도를 더한 행렬을 Cholesky 분해해 표집한다. 평활 분산 $\tau_k$는 역감마 조건부분포에서 뽑는다.

날짜 효과 $u_d$는 켤레 정규 조건부분포에서 뽑는다. 사전분포 $\mathcal N(0,\sigma^2_{u,g})$와 그날 백색화 잔차의 가중 우도를 결합한 것이다.

생산 이득의 로그 $\omega_{c,k}$와 자기상관 $\rho_g$는 랜덤워크 Metropolis로 갱신한다. $\omega_{c,k}$의 제안은 $(\mu_{k,\cdot},\gamma_k)$를 적분한 주변우도로 채택 여부를 정한다. $\rho_g$는 $[0,\,0.98)$ 밖의 제안을 기각한다. 혁신 분산 $\sigma^2_{\eta,g}$와 날짜 효과 분산 $\sigma^2_{u,g}$는 역감마 조건부분포에서 뽑는다.

표집은 6,000회 반복하고 앞 3,000회를 버린다. 남은 3,000개 사후표본이 예측 경로 3,000개가 된다. 수렴은 시드가 다른 사슬 사이에서 고전적 split-$\widehat R$ \citep{gelman1992rhat}로 점검하고, 1.05를 넘는 모수가 있으면 수렴 실패로 표시한다. 사전분포의 구체적인 값과 한 번의 반복 절차는 부록~\ref{app:model}에 둔다.

\subsection{예측 경로와 피크 위험}
\label{sec:ch2-paths}

예측분포는 사후표본마다 하루 경로 하나를 만들어 얻는다. 사후표본 $s=1,\dots,S$($S=3000$)의 모수로 예측일 계획에서 평균 곡선 $m^{(s)}_{d,t}$(식~\ref{eq:ch2-cbat}의 앞 세 항)를 계산하고 새 날짜 효과와 AR(1) 잡음을 더한다.
\begin{equation}
\tilde y^{(s)}_{d,t} = \max\!\Big\{0,\ s_y\big(m^{(s)}_{d,t} + u^{(s)}_d + e^{(s)}_{d,t}\big)\Big\},
\qquad u^{(s)}_d\sim\mathcal N\big(0,\sigma^{2(s)}_{u,g}\big),\quad
e^{(s)}_{d,t}=\rho^{(s)}_g e^{(s)}_{d,t-1}+\eta^{(s)}_{d,t}.
\label{eq:ch2-path}
\end{equation}
잡음 상태 $g$는 그날 계획의 가동유형으로 정한다. 전력은 음수가 될 수 없으므로 0에서 자른다. 난수 시드는 예측 날짜로 정해 같은 입력에서 같은 경로가 나온다.

슬롯 예측과 피크 분포는 모두 이 경로 묶음에서 계산한다. 점예측은 슬롯별 경로 중앙값이고, 분위수는 경로의 5\%에서 95\%까지 5\% 간격 분위수(q05--q95)다. 일 피크의 분포는 경로마다 구한 최댓값 $\tilde M^{(s)}_d=\max_t \tilde y^{(s)}_{d,t}$의 분포이며, 피크 점예측은 그 중앙값이다. 이는 슬롯별 중앙값 곡선의 최댓값 $\max_t\operatorname{median}_s(\tilde y^{(s)}_{d,t})$과 일반적으로 다르다. 피크 초과확률은 최댓값이 기준을 넘는 경로의 비율이다.
\begin{equation}
\hat P(M_d > C) = \frac1S\sum_{s=1}^{S}\mathbf 1\big[\tilde M^{(s)}_d > C\big],
\qquad C\in\{C_{50},C_{75},C_{90}\}.
\label{eq:ch2-risk}
\end{equation}
기준 $C_{50},C_{75},C_{90}$은 학습 구간 가동일 가운데 결측 슬롯이 4개 이하인 날의 관측 일 피크의 50·75·90\% 분위수다. 학습 구간이 바뀌면 기준도 다시 계산한다.

발행 시에는 96슬롯 전체에서 경로 최댓값을 구한다. 채점 시에는 관측 슬롯 $\mathcal O_d$에 대응하는 경로만으로 $\max_{t\in\mathcal O_d}\tilde y^{(s)}_{d,t}$를 계산해 관측 피크와 비교하고, 결측이 4슬롯을 넘는 날은 피크 평가에서 뺀다. 이 마스크는 사후 채점 규칙이며 발행 예측에 미래의 결측 여부를 입력하는 규칙이 아니다.

발행하는 위험확률은 식~\eqref{eq:ch2-risk}의 보정 전 경로 확률이다. 처음에는 예측일 직전 7일의 내부 구간에서 Platt 방식으로 한 번 더 보정했다. 그러나 이 짧은 보정 구간에는 $C_{90}$ 초과일이 없어 절편만 추정되었고 9월 가동일의 $C_{90}$ 위험이 모두 같은 값으로 고정되었다. 개발 구간에서도 $C_{90}$ AUC가 보정 전 0.926, 7일 Platt 0.788로 보정 전이 높았다. 그래서 9월 결과를 보기 전에 보정 전 확률을 발행하기로 사전 등록하였다.

\subsection{월 최대 갱신 경보 규칙}
\label{sec:ch2-alarm}

월 최대 갱신 경보는 그날 요금 적용 최대수요가 이미 확정된 기준선을 넘을 위험을 비용으로 환산해 경보를 낸다. 기본요금은 요금 적용 시간대(일요일·공휴일을 뺀 09--23시의 중간·최대부하 시간대)의 최대수요로 매겨진다. 이미 그 달과 래칫 기간에 기록된 최대보다 낮은 피크는 요금을 바꾸지 않는다. 따라서 위험한 날은 피크가 높은 날이 아니라 기준선을 새로 넘는 날이다.

기준선 $B_d$는 00:00에 알려진 관측만으로 정한다.
\begin{equation}
B_d = \max\Big\{\max_{j<d,\ j\in\text{같은 달}} M^{\mathrm{band}}_j,\ \ \max_{j<d,\ j\in\text{지난 1년의 래칫 적용월}} M^{\mathrm{band}}_j\Big\}.
\label{eq:ch2-baseline}
\end{equation}
$M^{\mathrm{band}}_j$는 $j$일 요금 적용 시간대의 관측 최대이고 래칫 적용월은 12·1·2·7·8·9월이다. 요금시간대가 없는 날은 경보 대상에서 제외한다. 과거의 유효 관측이 전혀 없으면 $B_d=0$이다. 자료가 2021년 1월에 시작해 2020년 12월 수요전력은 포함하지 못하므로, 이 기준선은 제공 자료로 관측 가능한 래칫이며 실제 청구 이력 전체를 재현한 값은 아니다.

C-BAT 경로에서 갱신 확률과 기대 초과 비용을 계산한다. 경로 $s$의 요금 적용 시간대 최대를 $\tilde M^{\mathrm{band},(s)}_d$라 하면,
\begin{equation}
P_d=\frac1S\sum_{s}\mathbf 1\big[\tilde M^{\mathrm{band},(s)}_d > B_d\big],
\qquad
E_d = 8{,}320\ \text{원/kW}\times\frac1S\sum_{s}\big(\tilde M^{\mathrm{band},(s)}_d - B_d\big)^{+}.
\label{eq:ch2-cost}
\end{equation}
8,320원/kW는 적용 요금제(선택 II)의 한 달 기본요금 환산단가다. $E_d$는 그날 새 최대가 생길 때 그 달 기본요금이 늘어나는 금액의 기댓값이다. 래칫으로 이후 달에 미치는 영향은 포함하지 않는다. 개발 구간 백테스트에서는 각 시드의 $P_d$와 $E_d$를 계산한 뒤 시드 0·1·2의 평균으로 경보 여부를 정한다.

경보는 기대 초과 비용이 점검 비용 이상일 때 낸다. 규칙은 $E_d\ge c$이며, 주 가정은 담당자 1--2시간의 점검 비용인 $c=30{,}000$원이다. 민감도 분석에는 $c=10{,}000$원, $100{,}000$원과 확률 기준 $P_d\ge0.5$를 쓴다. 이 규칙은 확률만 보는 경보와 달리 기준선을 얼마나 넘을지까지 반영한다. 경보를 받고 조치했을 때의 절감 효과는 이 규칙이 추정하지 않는다. 위험확률과 마찬가지로 Platt 보정 없이 경로 확률을 쓴다.
